\section{USAMO 1998}
        \begin{enumerate}
        	\item (\textit{USAMO 1998}) Suppose that the set $\{1,2,\cdots, 1998\}$ has been partitioned into disjoint pairs $\{a_i,b_i\}$ ($1\leq i\leq 999$) so that for all $i$, $|a_i-b_i|$ equals $1$ or $6$. Prove that the sum

 \[|a_1-b_1|+|a_2-b_2|+\cdots +|a_{999}-b_{999}|\]
ends in the digit $9$.


 



\textbf{Solution}

Notice that $|a_i - b_i| \equiv 1 \pmod 5$, so $S=|a_1-b_1|+|a_2-b_2|+\cdots +|a_{999}-b_{999}| \equiv 1+1+\cdots + 1 \equiv 999 \equiv 4 \bmod{5}$.

 Also, for integers $M, N$ we have $|M-N| \equiv M-N \equiv M+N \bmod{2}$.

 Thus, we also have $S \equiv a_1+b_1+a_2+b_2+\cdots +a_{999}+b_{999} \equiv 1+2+ \cdots +1998 \equiv 999*1999 \equiv 1 \bmod{2}$ also, so by the Chinese Remainder Theorem $S \equiv 9\bmod{10}$. Thus, $S$ ends in the digit 9, as desired.

 


	\item (\textit{USAMO 1998}) Let ${\cal C}_1$ and ${\cal C}_2$ be  concentric circles, with ${\cal C}_2$ in the interior of  ${\cal C}_1$. From a point $A$ on ${\cal C}_1$ one draws the tangent $AB$ to ${\cal C}_2$ ($B\in {\cal C}_2$). Let $C$ be the second point of intersection of $AB$ and ${\cal C}_1$, and let $D$ be the midpoint of $AB$. A line passing through $A$ intersects ${\cal C}_2$ at $E$ and $F$ in such a way that the perpendicular  bisectors of $DE$ and $CF$ intersect at a point $M$ on $AB$. Find, with proof,  the ratio $AM/MC$.


 



\textbf{Solution 1}


                    \begin{figure}[H]
                        \centering
                        \includegraphics[width=0.5\textwidth]{Solution/USAMO_I_1998_2_Solution_Figure_1.png}
                    \end{figure}
                    

 

 First, $AD=\frac{AB}{2}=\frac{AC}{4}$. Because $E$,$F$ and $B$ all lie on a circle, $AE \cdot AF=AB \cdot AB=\frac{AB}{2} \cdot 2AB=AD \cdot AC$. Therefore, $\triangle ACF \sim \triangle AED$, so $\angle ACF = \angle AED$. Thus, quadrilateral $CFED$ is cyclic, and $M$ must be the center of the circumcircle of $CFED$, which implies that $MC=\frac{CD}{2}$. Putting it all together,

 $\frac{AM}{MC}=\frac{AC-MC}{MC}=\frac{AC-\frac{CD}{2}}{\frac{CD}{2}}=\frac{AC-\frac{AC-AD}{2}}{\frac{AC-AD}{2}}=\frac{AC-\frac{3AC}{8}}{\frac{3AC}{8}}=\frac{\frac{5AC}{8}}{\frac{3AC}{8}}=\boxed{\frac{5}{3}}.$

 Borrowed from \href{https://mks.mff.cuni.cz/kalva/usa/usoln/usol982.html}{https://mks.mff.cuni.cz/kalva/usa/usoln/usol982.html}

 



\textbf{Solution 2}

We will use signed lengths. WLOG let $AC = 4$, $AM = x$, and $AE = y > 0$. Then $AB = 2$ and $AD = 1$. Therefore, $DM = x - 1$ and $MC = 4 - x$. By Power of a Point, $AE \cdot AF = AB^2 = 4$, so $AF = \frac{4}{AE} = \frac{4}{y}$, and $EF = \frac{4}{y} - y = \frac{4 - y^2}{y}$. Since $M$ is on the perpendicular bisectors of $DE$ and $CF$, we have $|DM| = |EM| = |x - 1|$ and $|CM| = |FM| = |4 - x|$.

 By Stewart's Theorem on $\triangle MFA$ and cevian $DM$, 
 \[\frac{4}{y}(y(\frac{4-y^2}{y}) + (x-1)^2) = y(4-x)^2 + \frac{4-y^2}{y} \cdot x^2;\]
 \[4(4 - y^2 + (x-1)^2) = y^2(4-x)^2 + (4-y^2)x^2;\]
 \[16 - 4y^2 + 4x^2 - 8x + 4 = x^2y^2 - 8xy^2 + 16y^2 + 4x^2 - x^2y^2;\]
 \[8xy^2 - 8x- 20y^2 + 20 = (8x - 20)(y^2 - 1) = 0.\]

 So either $8x - 20 = 0$ or $y^2 - 1 = 0$. If $y^2 - 1 = 0$, then $AE = y = 1 = AD$ and $AF = AC = 4$, so the perpendicular bisectors of $CF$ and $DE$ are the same line, and they do not intersect at a point. Therefore, $AM = x = \frac{5}{2}$ and $MC = \frac{3}{2}$, so $\frac{AM}{MC} = \boxed{\frac{5}{3}}$.

 


	\item (\textit{USAMO 1998}) Let $a_0,a_1,\cdots ,a_n$ be numbers from the interval $(0,\pi/2)$ such that

 \[\tan \left(a_0-\frac{\pi}{4}\right)+ \tan \left(a_1-\frac{\pi}{4}\right)+\cdots +\tan \left(a_n-\frac{\pi}{4}\right)\geq n-1.\]
Prove that

 \[\tan a_0\tan a_1 \cdots \tan a_n\geq n^{n+1}.\]


 



\textbf{Solution}

Let $y_i = \tan{\left(a_i - \dfrac{\pi}{4}\right)}$, where $0 \leqslant i \leqslant n$. Then we have:

 
                \begin{itemize}
                    \item $y_0 + y_1 + \cdots + y_n \geqslant n - 1$,\item $1 + y_i \geqslant \sum_{j \neq i}{(1 - y_j)}$,\item $\frac{1 + y_i}{n} \geqslant \frac{1}{n}\sum_{j \neq i}{(1 - y_j)}$.
                \end{itemize}


                By AM-GM:

 
 \begin{align*} \frac{1}{n}\sum_{j \neq i}{(1 - y_j)} &\geqslant \prod_{j \neq i}{(1 - y_j)^{\frac{1}{n}}} \\ \frac {1 + y_i}{n} &\geqslant \prod_{j \neq i}{(1 - y_j)^{\frac{1}{n}}} \\ \prod_{i = 0}^n\frac{1 + y_i}{n} &\geqslant \prod_{i = 0}^{n} {\prod_{j \neq i}{(1 - y_j)}^{\frac{1}{n}}} \\ \prod_{i = 0}^n\frac{1 + y_i}{n} &\geqslant \prod_{i = 0}^n{(1 - y_i)} \\ \prod_{i = 0}^n\frac{1 + y_i}{1 - y_i} &\geqslant \prod_{i = 0}^n{n} \\ \prod_{i = 0}^n\frac{1 + y_i}{1 - y_i} &\geqslant n^{n + 1}. \end{align*}


 Note that by the addition formula for tangents:

 
 \[\tan{(a_i)} = \tan{\left[\left(a_i - \frac{\pi}{4}\right) + \frac{\pi}{4}\right]} = \frac{1 + \tan{\left(a_i - \frac{\pi}{4}\right)}}{1 - \tan{\left(a_i - \frac{\pi}{4}\right)}} = \frac{1 + y_i}{1 - y_i}.\]

 So $\prod_{i = 0}^n{\frac{1 + y_i}{1 - y_i}} = \tan{(a_0)}\tan{(a_1)} \cdots \tan{(a_n)} \geqslant n^{n + 1}$, as desired. $\blacksquare$

 


	\item (\textit{USAMO 1998}) A computer screen shows a $98 \times 98$ chessboard, colored in the usual way. One can select with a mouse any rectangle with sides on the lines of the chessboard and click the mouse button: as a result,  the colors in the selected rectangle switch (black becomes white, white becomes black). Find, with proof, the minimum number of mouse clicks needed to make the chessboard all one color.


 



\textbf{Solution}

Answer: $98$.

 There are $4\cdot97$ adjacent pairs of squares in the border and each pair has one black and one white square. Each move can fix at most $4$ pairs, so we need at least $97$ moves. However, we start with two corners one color and two another, so at least one rectangle must include a corner square. But such a rectangle can only fix two pairs, so at least $98$ moves are needed.

 It is easy to see that 98 moves suffice: take 49 $1\times98$ rectangles (alternate rows), and 49 $98\times1$ rectangles (alternate columns).

 credit: \href{https://mks.mff.cuni.cz/kalva/usa/usoln/usol984.html}{https://mks.mff.cuni.cz/kalva/usa/usoln/usol984.html}

 editor: Brian Joseph

 second editor: integralarefun

 


	\item (\textit{USAMO 1998}) Prove that for each $n\geq 2$, there is a set $S$ of $n$ integers such that $(a-b)^2$ divides $ab$ for every distinct $a,b\in S$.


 



\textbf{Solution}

Proof by induction.  For n=2, the proof is trivial, since $S = (1,2)$ satisfies the condition.  Assume now that there is such a set S of n elements, $a_1, a_2,...a_n$ which satisfy the condition.  The key is to note that if $m=a_1a_2...a_n$, then if we define $b_i=a_i + km$ for all $i\le n$, where k is a positive integer, then $a_i \mid b_i$ and $b_i - b_j = a_i - a_j$, and so $(b_i - b_j)^2 = (a_i - a_j)^2 \mid a_ia_j \mid b_ib_j$.

 Let $b_{n+1}=m +km$.   Consider the set $T = (b_1,b_2,...,b_n,b_{n+1})$. To finish the proof, we simply need to choose a k such that $(b_{n+1}-b_i)^2 \mid b_{n+1}b_i$ for all $i\le n$.  Since $(b_{n+1}-b_i)^2 = (m-a_i)^2$, simply choose k so that $k+1 = (m-a_1)^2(m-a_2)^2...(m-a_n)^2$.

 


	\item (\textit{USAMO 1998}) Let $n \geq 5$ be an integer. Find the largest integer $k$ (as a function of $n$) such that there exists a convex $n$-gon $A_{1}A_{2}\dots A_{n}$ for which exactly $k$ of the quadrilaterals $A_{i}A_{i+1}A_{i+2}A_{i+3}$ have an inscribed circle. (Here $A_{n+j} = A_{j}$.)


 



\textbf{Solution}

Lemma: If quadrilaterals $A_iA_{i+1}A_{i+2}A_{i+3}$ and $A_{i+2}A_{i+3}A_{i+4}A_{i+5}$ in an equiangular $n$-gon are tangential, and $A_iA_{i+3}$ is the longest side quadrilateral $A_iA_{i+1}A_{i+2}A_{i+3}$ for all $i$, then quadrilateral $A_{i+1}A_{i+2}A_{i+3}A_{i+4}$ is not tangential.

 Proof:

 
                    \begin{figure}[H]
                        \centering
                        \includegraphics[width=0.5\textwidth]{Solution/USAMO_I_1998_6_Solution_Figure_1.png}
                    \end{figure}
                    

 If quadrilaterals $A_iA_{i+1}A_{i+2}A_{i+3}$ and $A_{i+2}A_{i+3}A_{i+4}A_{i+5}$ are tangential, then $A_iA_{i+3}$ must have side length of $j+y-u$, and $A_{i+2}A_{i+5}$ must have side length of $h + y - n$ (One can see this from what is known as walk-around). Suppose quadrilateral $A_{i+1}A_{i+2}A_{i+3}A_{i+4}$ is tangential. Then, again, we see that $A_{i+1}A_{i+4}$ must have side length $u + n - y$. We assumed by lemma that $A_iA_{i+3} > A_{i}A_{i+1}$ for all $i$, so we have $A_iA_{i+3} > j$, $A_{i+1}A_{i+4} > y$, and $A_{i+2}A_{i+5} > h$. If we add up the side lengths $A_iA_{i+3} + A_{i+1}A_{i+4} + A_{i+2}A_{i+5}$, we get:
 \[A_iA_{i+3} + A_{i+1}A_{i+4} + A_{i+2}A_{i+5} = j + y - u + h + y - n + u + n - y\]
 \[A_iA_{i+3} + A_{i+1}A_{i+4} + A_{i+2}A_{i+5} = j + h + y\]

 However, by the lemma, we assumed that $A_iA_{i+3} > j$, $A_{i+1}A_{i+4} > y$, and $A_{i+2}A_{i+5} > h$. Adding these up, we get:
 \[A_iA_{i+3} + A_{i+1}A_{i+4} + A_{i+2}A_{i+5} > j + h + y,\]

 which is a contradiction. Thus, quadrilateral $A_{i+1}A_{i+2}A_{i+3}A_{i+4}$ is not tangential, proving the lemma.

 


 By lemma, the maximum number of quadrilaterals in a $n$-gon occurs when the tangential quadrilaterals alternate, giving us $k = \lfloor \frac{n}{2} \rfloor$.

 
Note that one can find the ratio of side of an equiangular $n$-gon in order for alternating quadrilaterals to be tangential.

 
                    \begin{figure}[H]
                        \centering
                        \includegraphics[width=0.5\textwidth]{Solution/USAMO_I_1998_6_Solution_Figure_2.png}
                    \end{figure}
                    

 Since exterior angles of a equiangular $n$-gon have degree measure $\frac{2pi}{n}$, one can write the equation:

 
 \[2x = y + y + 2x \cos \frac{2pi}{n}\]
 \[y = x \left( 1- \cos \frac{2pi}{n} \right)\]
 \[\frac{y}{x} = \frac{1}{1- \cos \frac{2pi}{n}}\]

 Thus, we can find the ratio of sides of an equiangular $n$-gon which fits the maximum to be $1 : 1- \cos \frac{2\pi}{n}$. Note that if $n$ is even, we can easily alternate them, but if $n$ is odd, we must have two adjacent sides be the same length, and that length must be the larger side in the ratio  of adjacent sides. The proof is left as an exercise for the reader.

 


\end{enumerate}